\documentclass[10pt,a4paper]{amsart}
\usepackage{amsmath,amssymb,amsthm,graphicx}
\usepackage{verbatim}
\newenvironment{reference}{\comment}{\endcomment}
%\newenvironment{reference}{}{}
\newenvironment{slogan}{\comment}{\endcomment}
\newenvironment{history}{\comment}{\endcomment}
% For commutative diagrams we use Xy-pic
\usepackage[all]{xy}

% We use 2cell for 2-commutative diagrams.
\xyoption{2cell}
% Theorem environments.
%
\theoremstyle{plain}
\newtheorem{theorem}[subsection]{Theorem}
\newtheorem{proposition}[subsection]{Proposition}
\newtheorem{lemma}[subsection]{Lemma}

\theoremstyle{definition}
\newtheorem{definition}[subsection]{Definition}
\newtheorem{example}[subsection]{Example}
\newtheorem{exercise}[subsection]{Exercise}
\newtheorem{situation}[subsection]{Situation}

\theoremstyle{remark}
\newtheorem{remark}[subsection]{Remark}
\newtheorem{remarks}[subsection]{Remarks}

\numberwithin{equation}{subsection}

% Macros
%
\def\lim{\mathop{\mathrm{lim}}\nolimits}
\def\colim{\mathop{\mathrm{colim}}\nolimits}
\def\Spec{\mathop{\mathrm{Spec}}}
\def\Hom{\mathop{\mathrm{Hom}}\nolimits}
\def\Ext{\mathop{\mathrm{Ext}}\nolimits}
\def\SheafHom{\mathop{\mathcal{H}\!\mathit{om}}\nolimits}
\def\SheafExt{\mathop{\mathcal{E}\!\mathit{xt}}\nolimits}
\def\Sch{\mathit{Sch}}
\def\Mor{\mathop{\mathrm{Mor}}\nolimits}
\def\Ob{\mathop{\mathrm{Ob}}\nolimits}
\def\Sh{\mathop{\mathit{Sh}}\nolimits}
\def\NL{\mathop{N\!L}\nolimits}
\def\CH{\mathop{\mathrm{CH}}\nolimits}
\def\proetale{{pro\text{-}\acute{e}tale}}
\def\etale{{\acute{e}tale}}
\def\QCoh{\mathit{QCoh}}
\def\Ker{\mathop{\mathrm{Ker}}}
\def\Im{\mathop{\mathrm{Im}}}
\def\Coker{\mathop{\mathrm{Coker}}}
\def\Coim{\mathop{\mathrm{Coim}}}

% Boxtimes
%
\DeclareMathSymbol{\boxtimes}{\mathbin}{AMSa}{"02}

%
% Macros for moduli stacks/spaces
%
\def\QCohstack{\mathcal{QC}\!\mathit{oh}}
\def\Cohstack{\mathcal{C}\!\mathit{oh}}
\def\Spacesstack{\mathcal{S}\!\mathit{paces}}
\def\Quotfunctor{\mathrm{Quot}}
\def\Hilbfunctor{\mathrm{Hilb}}
\def\Curvesstack{\mathcal{C}\!\mathit{urves}}
\def\Polarizedstack{\mathcal{P}\!\mathit{olarized}}
\def\Complexesstack{\mathcal{C}\!\mathit{omplexes}}
% \Pic is the operator that assigns to X its picard group, usage \Pic(X)
% \Picardstack_{X/B} denotes the Picard stack of X over B
% \Picardfunctor_{X/B} denotes the Picard functor of X over B
\def\Pic{\mathop{\mathrm{Pic}}\nolimits}
\def\Picardstack{\mathcal{P}\!\mathit{ic}}
\def\Picardfunctor{\mathrm{Pic}}
\def\Deformationcategory{\mathcal{D}\!\mathit{ef}}
\usepackage{cleveref}
\crefname{originallemma}{original lemma}{original lemmas}
\Crefname{originallemma}{Original lemma}{Original lemmas}
\makeatletter
\newlabel{lemma-exactness-pushforward-pullback}{{lemma-exactness-pushforward-pullback}{source}{}{sourcepointer-1}{}}
\newlabel{lemma-exactness-pushforward-pullback@cref}{{[originallemma][0][]lemma-exactness-pushforward-pullback}{[1][0][]source}{}{}{}}
\makeatother
\crefname{originalreference}{original reference}{original references}
\Crefname{originalreference}{Original reference}{Original references}
\makeatletter
\newlabel{sheaves-section-ringed-spaces-functoriality-modules}{{sheaves-section-ringed-spaces-functoriality-modules}{source}{}{sourcepointer-2}{}}
\newlabel{sheaves-section-ringed-spaces-functoriality-modules@cref}{{[originalreference][0][]sheaves-section-ringed-spaces-functoriality-modules}{[1][0][]source}{}{}{}}
\makeatother
\crefname{originalreference}{original reference}{original references}
\Crefname{originalreference}{Original reference}{Original references}
\makeatletter
\newlabel{sheaves-lemma-stalk-pullback-modules}{{sheaves-lemma-stalk-pullback-modules}{source}{}{sourcepointer-3}{}}
\newlabel{sheaves-lemma-stalk-pullback-modules@cref}{{[originalreference][0][]sheaves-lemma-stalk-pullback-modules}{[1][0][]source}{}{}{}}
\makeatother
\title{Three constructions of the quasi-coherent sheaf associated to a module over the ring of global sections agree}
\author{The Stacks Project Authors}
\date{stacks-project:modules.tex:local-snapshot-20261012; GNU FDL 1.2 or later}
\begin{document}
\maketitle
\noindent\textit{Editorial note.} All delivered mathematical text is the author's own, copied from the pinned Stacks Project chapter source sources/stacks-modules/pinned-original-source.tex; no statement or proof step was rewritten, reordered or removed.\par\medskip
\noindent\textit{Editorial note.} Preamble provenance: the Stacks Project keeps one single shared preamble for every chapter; this card takes it directly from the staged sources/stacks-modules/preamble.tex as three verbatim spans (lines 8-12, 14-18 and 154-217), so no derived file was written and nothing was added or redefined. The wrapper class and the base packages are supplied by the assembler.\par\medskip
\noindent\textit{Editorial note.} Bibliography provenance: neither the counted statement nor its proof contains a citation (this lane rescans the whole span before writing the card), so the card carries no bibliography and no entry was reconstructed.\par\medskip
\noindent\textit{Editorial note.} Selection provenance: this card was opened under the file-based control contract (CONTROL.md), with the mechanical pre-card gate run on its statement and proof interval before the card was written.\par\medskip
\noindent\textit{Allowed standard prerequisites (editorial).}\begin{itemize}
\item The pullback of modules along the structure morphism to the point with the given ring, and its description as a sheafification of a tensor product presheaf. Invoked inside the counted proof (source lines 1402-1405) as a cross-chapter pointer. Identifies the first and the third construction.
\item The right exactness of the pullback and its compatibility with direct sums. Invoked inside the counted proof (source lines 1406-1411) as a pointer to a lemma of this chapter. Transports a presentation of the module to a presentation of the pulled back sheaf, which identifies the second construction with the first.
\end{itemize}\medskip
\ifdefined\hypertarget\hypertarget{sourcepointer-1}{}\fi
\noindent\textit{Source reference.} Original lemma lemma-exactness-pushforward-pullback: sources/stacks-modules/pinned-original-source.tex, source lines 990--990. Navigation pointer only: the counted text refers to the Stacks item lemma-exactness-pushforward-pullback, which belongs to another chapter of the project and is not part of this card. The pointer carries the target's own label as its number because this chapter's staged snapshot does not contain the other chapter's tag table; its content is not reproduced here.\par\medskip
\ifdefined\hypertarget\hypertarget{sourcepointer-2}{}\fi
\noindent\textit{Source reference.} Original reference sheaves-section-ringed-spaces-functoriality-modules: sources/stacks-modules/pinned-original-source.tex, source lines 1405--1405. Navigation pointer only: the counted text refers to the Stacks item sheaves-section-ringed-spaces-functoriality-modules, which belongs to another chapter of the project and is not part of this card. The pointer carries the target's own label as its number because this chapter's staged snapshot does not contain the other chapter's tag table; its content is not reproduced here.\par\medskip
\ifdefined\hypertarget\hypertarget{sourcepointer-3}{}\fi
\noindent\textit{Source reference.} Original reference sheaves-lemma-stalk-pullback-modules: sources/stacks-modules/pinned-original-source.tex, source lines 1418--1418. Navigation pointer only: the counted text refers to the Stacks item sheaves-lemma-stalk-pullback-modules, which belongs to another chapter of the project and is not part of this card. The pointer carries the target's own label as its number because this chapter's staged snapshot does not contain the other chapter's tag table; its content is not reproduced here.\par\medskip
\begin{lemma}
\label{lemma-construct-quasi-coherent-sheaves}
Let $(X, \mathcal{O}_X)$ be ringed space.
Let $\alpha : R \to \Gamma(X, \mathcal{O}_X)$ be a ring homomorphism from
a ring $R$ into the ring of global sections on $X$.
Let $M$ be an $R$-module.
The following three constructions give canonically isomorphic
sheaves of $\mathcal{O}_X$-modules:
\begin{enumerate}
\item Let $\pi : (X, \mathcal{O}_X) \longrightarrow (\{*\}, R)$
be the morphism of ringed spaces with $\pi : X \to \{*\}$
the unique map and with $\pi$-map $\pi^\sharp$ the given map
$\alpha : R \to \Gamma(X, \mathcal{O}_X)$. Set $\mathcal{F}_1 = \pi^*M$.
\item Choose a presentation
$\bigoplus_{j \in J} R \to \bigoplus_{i \in I} R \to M \to 0$.
Set
$$
\mathcal{F}_2 = \Coker\left(
\bigoplus\nolimits_{j \in J} \mathcal{O}_X
\to
\bigoplus\nolimits_{i \in I} \mathcal{O}_X
\right).
$$
Here the map on the component $\mathcal{O}_X$ corresponding to $j \in J$
given by the section $\sum_i \alpha(r_{ij})$ where the $r_{ij}$
are the matrix coefficients of the map in the presentation of $M$.
\item Set $\mathcal{F}_3$ equal to the sheaf associated to the presheaf
$U \mapsto \mathcal{O}_X(U) \otimes_R M$, where the map
$R \to \mathcal{O}_X(U)$ is the composition of $\alpha$ and
the restriction map $\mathcal{O}_X(X) \to \mathcal{O}_X(U)$.
\end{enumerate}
This construction has the following properties:
\begin{enumerate}
\item The resulting sheaf of $\mathcal{O}_X$-modules
$\mathcal{F}_M = \mathcal{F}_1 = \mathcal{F}_2 = \mathcal{F}_3$
is quasi-coherent.
\item The construction gives a functor from
the category of $R$-modules to the category of quasi-coherent
sheaves on $X$ which commutes with arbitrary colimits.
\item For any $x \in X$ we have
$\mathcal{F}_{M, x} = \mathcal{O}_{X, x} \otimes_R M$
functorial in $M$.
\item Given any $\mathcal{O}_X$-module
$\mathcal{G}$ we have
$$
\Mor_{\mathcal{O}_X}(\mathcal{F}_M, \mathcal{G})
=
\Hom_R(M, \Gamma(X, \mathcal{G}))
$$
where the $R$-module structure on $\Gamma(X, \mathcal{G})$
comes from the $\Gamma(X, \mathcal{O}_X)$-module structure via
$\alpha$.
\end{enumerate}
\end{lemma}

\begin{proof}
The isomorphism between $\mathcal{F}_1$ and $\mathcal{F}_3$
comes from the fact that $\pi^*$ is defined as the sheafification
of the presheaf in (3), see Sheaves, Section
\ref{sheaves-section-ringed-spaces-functoriality-modules}.
The isomorphism between the constructions in (2) and (1) comes
from the fact that the functor $\pi^*$ is right exact, so
$\pi^*(\bigoplus_{j \in J} R) \to \pi^*(\bigoplus_{i \in I} R) \to
\pi^*M \to 0$ is exact, $\pi^*$ commutes with arbitrary
direct sums, see Lemma \ref{lemma-exactness-pushforward-pullback},
and finally the fact that $\pi^*(R) = \mathcal{O}_X$.

\medskip\noindent
Assertion (1) is clear from construction (2).
Assertion (2) is clear since $\pi^*$ has these properties.
Assertion (3) follows from the description of stalks of
pullback sheaves, see
Sheaves, Lemma \ref{sheaves-lemma-stalk-pullback-modules}.
Assertion (4) follows from adjointness of $\pi_*$ and
$\pi^*$.
\end{proof}

\end{document}
